\section{Metodologia de Coleta e Processamento}

Esta seção descreve o processo utilizado para construção do dataset proposto, incluindo as fontes de dados empregadas, as etapas de extração, transformação e integração das informações e a arquitetura adotada para armazenamento e disponibilização dos dados. O processo foi concebido para garantir reprodutibilidade, rastreabilidade e facilidade de atualização do conjunto de dados.

\subsection{Seleção das Fontes de Dados}

A seleção das fontes de dados foi conduzida por meio de uma busca exploratória em mecanismos de pesquisa, repositórios públicos e documentações oficiais relacionadas à segurança cibernética e vulnerabilidades de software. O objetivo foi identificar bases públicas capazes de fornecer informações complementares sobre identificação de vulnerabilidades, severidade, exploração conhecida, fraquezas de software e pacotes afetados, passíveis de integração por meio de identificadores compartilhados.

As fontes candidatas foram avaliadas quanto à disponibilidade pública dos dados, documentação técnica, forma de acesso (APIs, feeds, arquivos ou repositórios Git), frequência de atualização, granularidade das informações, restrições de licenciamento e potencial de reutilização em pesquisas acadêmicas. Foram consideradas apenas bases que disponibilizassem acesso aberto aos metadados e permitissem integração com outras fontes utilizando identificadores padronizados, principalmente o CVE-ID.

Ao final do processo, foram identificadas inicialmente 30 bases públicas, das quais 10 foram analisadas tecnicamente e 7 selecionadas para análise comparativa. A versão atual do datalake implementa quatro dessas fontes: CISA Known Exploited Vulnerabilities (KEV), CVE List V5, National Vulnerability Database (NVD) e GitHub Advisory Database. Além dessas, o modelo de dados prevê a futura incorporação de bases complementares, como CWE, OSV, EPSS, Exploit-DB e outras fontes de threat intelligence.

As principais fontes utilizadas na implementação são apresentadas na Tabela~\ref{tab:data_sources}.

\begin{table}[htbp]
    \centering
    \caption{Fontes implementadas na versão atual do dataset}
    \label{tab:data_sources}
    \small
    \renewcommand{\arraystretch}{1.15}
    \begin{tabular}{|>{\raggedright\arraybackslash}p{3.2cm}|>{\raggedright\arraybackslash}p{2.8cm}|>{\raggedright\arraybackslash}p{6.5cm}|}
        \hline
        \textbf{Fonte}           &
        \textbf{Forma de acesso} &
        \textbf{Principais informações}                                                \\
        \hline

        CISA KEV                 &
        JSON Feed                &
        Exploração conhecida, data de inclusão, ação recomendada e prazo de mitigação. \\

        \hline

        CVE List V5              &
        Repositório Git          &
        Registro canônico de CVEs, descrições, produtos afetados e referências.        \\

        \hline

        NVD                      &
        REST API                 &
        CVSS, CWE, CPE, métricas de severidade e metadados técnicos.                   \\

        \hline

        GitHub Advisory Database &
        Repositório Git          &
        Pacotes afetados, ecossistemas, versões vulneráveis/corrigidas e advisories.   \\

        \hline
    \end{tabular}
\end{table}


\subsection{Arquitetura}
\label{subsec:arquitetura}

A arquitetura do pipeline segue o padrão de camadas de maturidade de dados (\textit{medallion architecture}), organizando o fluxo em três estágios sucessivos, \textbf{bronze} (dados brutos), \textbf{silver} (dados processados e integrados) e \textbf{gold} (dados curados em grafo), cada um implementado como um estágio isolado do pipeline (extração, transformação e carga, respectivamente), conforme detalhado na Seção~\ref{subsec:pipeline}. A opção por essa arquitetura busca favorecer a rastreabilidade (é sempre possível recuperar o dado original na camada bronze), a modularidade (cada camada é produzida por um conjunto independente de scripts) e a reprodutibilidade do processo como um todo.

O repositório do projeto reflete essa separação em sua estrutura de diretórios: \texttt{scripts/} concentra os módulos de extração, transformação, carga e análise; \texttt{datalake/} armazena fisicamente os dados segundo seu grau de maturidade; \texttt{config/} reúne parâmetros e variáveis globais; e \texttt{docs/} contém documentação técnica complementar.

Do ponto de vista tecnológico, o pipeline foi implementado em Python 3.12+, com a biblioteca \textit{Requests} para consumo de APIs externas e arquivos YAML/dotenv para gerenciamento de configurações. O processamento é realizado com o DuckDB, motor SQL de alto desempenho para execução \textit{in-process}, enquanto a orquestração das etapas é feita via Apache Airflow, containerizado com Docker --- o que garante agendamento automatizado, reprodutibilidade e monitoramento centralizado das execuções, com logs persistidos tanto na interface do Airflow quanto em arquivos locais.

\subsection{Pipeline ETL/ELT}
\label{subsec:pipeline}

O pipeline foi implementado integralmente em Python, utilizando o DuckDB como motor analítico \textit{in-memory} para as etapas de transformação e junção. A arquitetura segue um modelo em camadas: \textit{bronze} (dados brutos), \textit{silver} (dados integrados e limpos) e \textit{gold} (materialização em grafo de propriedades), organizadas no diretório \texttt{datalake/}.


\subsubsection{Camada Bronze -- Extração}

Cada fonte possui um script dedicado de extração. Para a \textbf{CISA KEV}, a coleta realiza uma requisição HTTP ao endpoint JSON oficial da agência, com política de retentativa exponencial via biblioteca \texttt{tenacity}. Para a \textbf{CVE List V5}, o script baixa o repositório oficial do CVEProject no GitHub em formato ZIP, percorre recursivamente os arquivos JSON e gera um arquivo JSONL consolidado.

A extração do \textbf{NVD} utiliza a API REST 2.0 com paginação, controle de taxa, retentativa automática e suporte à coleta incremental. O processamento do \textbf{GitHub Advisories} realiza o \textit{flattening} das estruturas aninhadas e gera um conjunto tabular. A coleta do \textbf{CWE} é realizada a partir do portal oficial da MITRE, preservando os arquivos CSV originais para fins de rastreabilidade.

\subsubsection{Camada Silver -- Transformação}

A camada \textit{silver} realiza a limpeza, normalização e integração dos dados utilizando DuckDB. O processamento do CWE remove duplicidades, extrai relacionamentos entre fraquezas e gera relatórios de qualidade. O processamento das CVEs utiliza o NVD como fonte âncora, enriquecendo os registros com informações provenientes da CVE List V5, CISA KEV e GitHub Advisories.

As junções são realizadas por \texttt{LEFT JOIN} utilizando o identificador CVE. Após a integração, são aplicadas etapas de padronização dos atributos, conversão de tipos, preenchimento de valores ausentes e gravação do resultado em \texttt{Parquet} com compressão Snappy.

\subsubsection{Camada Gold -- Carga}

A camada \textit{gold} é responsável pela carga dos dados no Neo4j. O processo cria ou atualiza os nós \texttt{:CVE}, \texttt{:CWE}, \texttt{:Package} e \texttt{:Reference}, além dos relacionamentos \texttt{HAS\_WEAKNESS}, \texttt{AFFECTS\_PACKAGE}, \texttt{HAS\_REFERENCE} e dos relacionamentos dinâmicos entre CWEs.

São criados índices sobre os principais atributos de busca (\texttt{CVE.id}, \texttt{CWE.id}, \texttt{CVE.severity}, \texttt{Package.ecosystem} e \texttt{Reference.url}), garantindo melhor desempenho das consultas. A carga é idempotente, permitindo reexecuções sem duplicação de nós ou relacionamentos.



